\begin{annexitem}{annex:p19.rat46.fixed-event}{thm:p19-rat46-fixed-event}{The complete finite truth-table obstruction}
\begin{proof}
Write the fixed event's membership coefficients on
\(A C,A\neg C,\neg A C,\neg A\neg C\) as
\(d_{11},d_{10},d_{01},d_{00}\), each either zero or one.
For parameters \(s,r,q\in(0,1)\), assign the four atoms masses
\(sr,s(1-r),(1-s)q,(1-s)(1-q)\). All are positive, sum to one,
and give \(P(A)=s\) and \(P(C\mid A)=r\).
The fixed event has probability
\[
 d_{11}sr+d_{10}s(1-r)
       +d_{01}(1-s)q+d_{00}(1-s)(1-q).
\]
If this equaled \(r\) for every choice, changing \(q\) with
\(s,r\) fixed would force \(d_{01}=d_{00}\).
Now choose two distinct \(r\)'s with \(s,q\) fixed and subtract
the claimed equalities. Cancellation gives
\(s(d_{11}-d_{10})=1\). But
\(d_{11}-d_{10}\in\{-1,0,1\}\) and \(0<s<1\), so that is
impossible. Thus no such fixed event exists. The proof requires no
limiting zero masses and does not apply to an assessment procedure,
a probability-dependent event selection or a conditional random quantity
outside this Boolean-event type.
\end{proof}
\end{annexitem}

\begin{annexitem}{annex:p19.rat46.frequency-contrast}{thm:p19-rat46-frequency-contrast}{Every denominator in the two source illustrations}
\begin{proof}
For a uniform draw from the 100-card illustration,
\(P(A)=60/100\), \(P(A\land C)=50/100\), and
\(P(A\land\neg C)=10/100\). The positive-antecedent ratio is
\((50/100)/(60/100)=5/6\); the conjunction is \(1/2\);
and material implication is the complement of the falsifying case,
\(1-10/100=9/10\). The other two counts make the total 100 and
do not enter the ratio's denominator.

For the six equiprobable die faces, yellow occupies four faces and
yellow-and-circle two. Thus the ratio is \((2/6)/(4/6)=1/2\),
the conjunction is \(2/6=1/3\), and the material falsifier is the
two yellow-square faces, giving \(1-2/6=2/3\). These calculations
compare three specified objects; they do not choose which object a
participant intended, and do not calculate the probability that a rule
holds for every card in an experimentally sampled ten-card collection.
\end{proof}
\end{annexitem}

\begin{annexitem}{annex:p19.rat46.model-parameters}{thm:p19-rat46-model-parameters}{Joint masses, feasibility and the MP--MT qualification}
\begin{proof}
By the definition of \(e\), \(P(A\land C)=a(1-e)\).
Subtract this from \(P(A)=a\) to obtain \(P(A\land\neg C)=ae\).
Subtract it from \(P(C)=b\) to obtain \(P(\neg A\land C)=b-a(1-e)\).
The fourth mass is one minus the first three, namely \(1-b-ae\).
Their sum is one; nonnegativity is exactly the bounds stated in the
definition, together with \(0\leq e\leq1\) and the stipulated
\(0<a<1\). Thus this is a well-defined joint probability distribution
under those conditions.

Divide the first mass by \(a\) for MP, the fourth by \(1-b\)
for MT, the first by \(b\) for AC, and the fourth by \(1-a\)
for DA. This gives the four expressions with precisely the stated
positive-denominator restrictions. A zero denominator is not resolved
by assigning the displayed fraction a value.

For \(a=0.5,e=0.01,b=0.994\), the masses in the same order are
\(0.495,0.005,0.499,0.001\). They are strictly positive and sum to
one. MP is \(0.495/0.5=0.99\), whereas MT is
\(0.001/0.006=1/6\). Generally the MT error is
\(ae/(1-b)\), not simply \(e\). If \(a/(1-b)\) is bounded
while \(e\) tends to zero, MT tends to one; without that condition
the counterexample demonstrates the missing inference. This repairs
the formal implication while leaving the separately attributed fitted-
model criticism to its empirical evidence.
\end{proof}
\end{annexitem}
